\documentclass[10pt,letterpaper]{article}
\usepackage[top=0.85in,left=2.75in,footskip=0.75in]{geometry}

% amsmath and amssymb packages, useful for mathematical formulas and symbols
\usepackage{amsmath,amssymb}

% Use adjustwidth environment to exceed column width (see example table in text)
\usepackage{changepage}

% Use Unicode characters when possible
\usepackage[utf8x]{inputenc}

% textcomp package and marvosym package for additional characters
\usepackage{textcomp,marvosym}

% cite package, to clean up citations in the main text. Do not remove.
\usepackage{cite}

% Use nameref to cite supporting information files (see Supporting Information section for more info)
\usepackage{nameref,hyperref}

% line numbers
\usepackage[right]{lineno}

% ligatures disabled
\usepackage{microtype}
\DisableLigatures[f]{encoding = *, family = * }

% color can be used to apply background shading to table cells only
\usepackage[table]{xcolor}

% array package and thick rules for tables
\usepackage{array}

% enumerate package lets us use letters instead of numbers
\usepackage{enumerate}

% create "+" rule type for thick vertical lines
\newcolumntype{+}{!{\vrule width 2pt}}

% create \thickcline for thick horizontal lines of variable length
\newlength\savedwidth
\newcommand\thickcline[1]{%
  \noalign{\global\savedwidth\arrayrulewidth\global\arrayrulewidth 2pt}%
  \cline{#1}%
  \noalign{\vskip\arrayrulewidth}%
  \noalign{\global\arrayrulewidth\savedwidth}%
}

% \thickhline command for thick horizontal lines that span the table
\newcommand\thickhline{\noalign{\global\savedwidth\arrayrulewidth\global\arrayrulewidth 2pt}%
\hline
\noalign{\global\arrayrulewidth\savedwidth}}

\usepackage{color}

% Remove comment for double spacing
%\usepackage{setspace}
%\doublespacing

% Text layout
\raggedright
\setlength{\parindent}{0.5cm}
\textwidth 5.25in
\textheight 8.75in

% Bold the 'Figure #' in the caption and separate it from the title/caption with a period
% Captions will be left justified
\usepackage[aboveskip=1pt,labelfont=bf,labelsep=period,justification=raggedright,singlelinecheck=off]{caption}
\renewcommand{\figurename}{Fig}

% Use the PLoS provided BiBTeX style
\bibliographystyle{plos2015}

% Remove brackets from numbering in List of References
\makeatletter
\renewcommand{\@biblabel}[1]{\quad#1.}
\makeatother

% Leave date blank
\date{}

% Header and Footer with logo
\usepackage{lastpage,fancyhdr,graphicx}
\usepackage{epstopdf}
\pagestyle{myheadings}
\pagestyle{fancy}
\fancyhf{}
\setlength{\headheight}{27.023pt}
\lhead{\includegraphics[width=2.0in]{PLOS-submission.eps}}
\rfoot{\thepage/\pageref{LastPage}}
\renewcommand{\footrule}{\hrule height 2pt \vspace{2mm}}
\fancyheadoffset[L]{2.25in}
\fancyfootoffset[L]{2.25in}
\lfoot{\sf PLOS}


\newcommand{\rulemajor}[1]{\section*{#1}}
\newcommand{\withurl}[2]{{#1}\footnote{{\texttt{#2}}}}

\begin{document}
\vspace*{0.2in}

\begin{flushleft}
{\Large
\textbf\newline{Ten Quick Tips for Staying Safe Online}
}
\newline
\\
{Greg Wilson}\textsuperscript{1{\ddag}*}
\\
\bigskip
\textbf{1} RStudio, Inc., Toronto, ON, Canada\\
* Corresponding author, greg.wilson@rstudio.com. \\
\bigskip
{\ddag} These authors contributed equally to this work.
\end{flushleft}

\section*{Introduction}

Researchers studying everything from sexual health to COVID-19 to gun violence
are increasingly likely to be targeted because of their work.
While research institutions have rules and guidelines for safeguarding sensitive information,
these usually do not address the problem of keeping \emph{individuals} safe
from either targeted attacks (e.g., Climategate \cite{Natu2010})
or the kinds of ``drive-by'' threats that everyone now faces regardless of their occupation.

Hollywood depictions of everyday threats are as far from reality as their portrayals of scientists,
but more realistic guidance for personal digital security is now freely available \cite{FLD,EFJ2015,EFF}.
The ten quick tips in this paper are a starting point:
while they apply to everyone,
they are particularly relevant to researchers.
They do not guarantee complete safety,
any more than seatbelts guarantee safe driving,
but following them greatly reduces the likelihood of harm.

\rulemajor{Rule 1: Put on your own mask.}

The first and most important rule is that
we should not rely on companies, universities, and other institutions to protect us,
for the simple reason that they are not penalized if they don't.
As recently as ten years ago we could blame the lack of meaningful institutional liability for data breaches
on the law being slow to catch up with rapidly-changing technology.
Today,
it's clear that the law hasn't caught up because large tech companies don't want it to.
Weekly reports of data breaches have minimal impact on profitability and no individual ever goes to jail

Much of what institutions force us to go through online is \textbf{security theater}
intended to make us believe something is being done
rather than to actually make us safer.
Requiring people to take off their shoes at airports is perhaps the best-known example;
random searches of backpacks and bags at the entrance to the subway is another,
since it's hard to imagine that a would-be attacker \emph{wouldn't} just to another entrance.

Security theater is counter-productive because it encourages us to cut corners
in ways that actually make us \emph{less} safe.
For example,
forcing people to change passwords every three months
encourages people to choose memorable (and therefore easy-to-guess) passwords.

\rulemajor{Rule 2: Digital security is rarely the weakest link.}

The second rule is to remember that most attacks take place offline,
and that the most effective ones are often the simplest.
At an airport several years ago,
the author heard a professor of computer science try to re-set an online account over the phone.
In just a couple of minutes,
they had inadvertently told everyone in the lounge their full name,
their date of birth,
the three-digit verification code on the back of their credit card,
and what was almost certainly their mother's maiden name.

On that same trip,
a complete stranger (also a computer science professor)
asked the author to watch their laptop while they went to the washroom.
They were logged in to several online accounts at the time,
and even if they weren't,
the passwords saved in their browser could have been viewed with just a few mouse clicks.

The moral of these stories is that safety comes from good habits,
not technology.
\textbf{Social engineering} is far more common than hacking:
in practice it is far easier to trick someone into giving you their password
than to break into their devices digitally.

The key practice is \textbf{situational awareness},
which is a fancy way of saying,
``Pay attention to what's happening and respond accordingly.''
You shouldn't always be afraid to walk after dark,
but you should use good judgment about where not to go.
Similarly,
if you notice unusual activity or if you start working on something that may attract unwanted attention,
you should take more precautions than usual.
For example,
Boris should recognize that agreeing to be an expert witness
increases the odds that he will be targeted,
and should be more careful about what he puts into email while preparing and delivering his testimony.

The corollary to situational awareness is to de-escalate when you can.
Being on guard all the time is exhausting
and quickly leads to \textbf{security fatigue} \cite{Stan2016}.
If you are too tired to follow good practices,
knowing them does you no good.

\rulemajor{Rule 3: Use relevant threat models.}

Edward Snowden and the journalists who worked with him took extraordinary measures
to safeguard themselves against \textbf{state-level actors} \cite{Snow2019},
but most of us aren't involved in issues of national security
and don't need to take those kinds of precautions.
Instead,
we typically face one of three kinds of threat illustrated by the examples below.

\begin{itemize}

\item
  \textbf{Casual threats} are opportunistic.
  For example,
  Monica, a professor in psychology, is targeted by Mohan,
  an undergraduate in computer science
  who spends hours every day in online echo chambers complaining about how ``SJW bullshit'' is ruining tech.
  He really didn't enjoy Monica's guest lecture on discrimination and inclusivity in his software engineering class,
  and thinks it would be a laugh to make her the target of anonymous abuse online.
  He is unlikely to invest significant effort in his attack (at least not initially),
  but his attack may be backed up by more knowledgeable advisors in online forums.

\item
  \textbf{Intimate threats} come from people who know their targets' passwords
  or have a chance to install spyware on their targets' devices \cite{Leit2019}.
  For example,
  Elena, graduate student,
  is targeted by her romantic partner Eric,
  who is also a graduate student but not in the same department.
  Their relationship had become increasingly abusive over the last two years.
  With the help of friends,
  Elena has moved out of their shared apartment and is rebuilding her life;
  Eric is obsessed with the idea that she left him for someone else and is now stalking her.

\item
  \textbf{Insider threats} come from people who have legitimate access to accounts and devices.
  For example,
  Boris, professor of medicine,
  is targeted by Bethany, who works for the university's IT department.
  Boris has agreed to serve as an expert witness in an upcoming liability case involving a large chemical company;
  Bethany has been asked by a former colleague to find out what he is going to say in order to discredit his testimony.

\end{itemize}

\rulemajor{Rule 4: Use a password manager.}

Using a weak password is a good way to ensure that your account will eventually be compromised,
in part because \textbf{dictionary attacks} can be run offline against encrypted password files
to find passwords that match common patterns.
Using a clever password scheme,
such as the name of the site plus a word only you know,
does not increase security by much:
whatever scheme you have thought of,
attackers have seen before.
And since people are often identified on multiple sites by the same email address,
as soon as one site where you've used that scheme is compromised,
attackers can guess the scheme and use it elsewhere.

Reusing passwords ensures that damage spreads,
so using a different password for each site helps limit harm if any are compromised.
However,
strong passwords are hard to remember and to type,
so always use a \textbf{password manager}
that generates strong passwords
and saves them all under a master \textbf{passphrase}.
Your passphrase should be several words long
and something you are unlikely to forget.
It does create a single point of attack,
but is still safer than choosing passwords yourself,
since password managers aren't fooled by similar-seeming sites like \texttt{paypaI.com}.

And despite what you may have heard,
writing passwords down and keeping them in your wallet is not a bad practice:
you have been keeping bits of paper safe since you were a child,
and you'll know if your written passwords go missing or are stolen.

\rulemajor{Rule 5: Use two-factor authentication.}

Authentication relies on something you \emph{know} (like a password),
something you \emph{have} (like a security key),
or something you \emph{are} (like your fingerprints).
\textbf{Two-factor authentication} requires two of these together to establish your identity,
e.g.,
a password (which can be stolen electronically)
plus a random code generated by an app on your phone
(which means attackers need access to you).

2FA is as important to security as using a password manager,
but where possible,
you should rely on an app for 2FA instead of using text messages.
What you should \emph{never} do is share a confirmation code,
since a common attack is to trigger a password reset
and then call the victim pretending to be from the IT department
and ask them to read the code back to ``verify'' your account.
As soon as you do this,
the attacker can change your password and get into your account.

\begin{quote}
  Many security experts now recommend using a physical 2FA key such as a YubiKey,
  which fits on a keychain and plugs into a standard USB port.
  Sites like \href{https://techsolidarity.org/}{Tech Solidarity} have easy-to-follow instructions
  explaining how to set them up for a range of popular social networking sites.
\end{quote}

\rulemajor{Rule 6: Think before opening.}

Much of the software we use was designed in more innocent times,
and since companies are almost never held liable for the damage caused by their software,
they have consistently prioritized convenience for the many over harm to the few.
One common example is documents that contain code called ``macros''
that automatically execute when the document is opened.
Used for good,
a macro can check that an address field has been filled in correctly.
Used for evil,
it can mail the contents of your address book anywhere in the world.
Microsoft Word and Excel are particularly notorious for this vulnerability,
but many other kinds of documents have the same flaw.

Attempts to get you to open an email attachment,
click on a link,
install software,
or log into a website are called \textbf{phishing} attacks.
The strongest defense is to never do these things,
but in the modern world that would make most work impossible.
The second-best defense is to take sensible precautions.
For example,
\emph{never} open email attachments without first running them through a virus scanner---not even
if you trust the person who sent it,
because their computer might just have been compromised.
Using cloud applications such as Google Docs instead of desktop applications means you don't have to download files,
but the tradeoff is that you are now giving all of your data to a company
whose business model is to sell your personal information to unaccountable third parties.

Similarly,
don't click links in emails without checking them first:
instead,
hover over the link and see if it matches the site it claims to be.
Alternatively,
log into the site manually rather than following the provided link.
It takes more time,
but is still faster than fixing your credit rating.
And when you go to a web site,
check the real domain name in the URL:
\texttt{paypaI.com} with an upper-case ``I'' instead of a lower-case ``l'' is not the site it pretends to be,
and \texttt{wwwpaypal.com} is a different domain than \texttt{www.paypal.com}.

\begin{quote}
  Many sites send an email with a random URL to confirm your identity when you are resetting your password.
  On the one hand,
  this means that an attacker has to get access to your email in order to break into your account.
  On the other hand,
  random URLs are hard to type in,
  so these emails encourage us to click on links in emails.
  If you are not expecting a password reset email,
  \emph{don't click on the link}.
\end{quote}

While phishing attacks are wide-ranging,
\textbf{spearphishing} uses data harvested from previous victims to attack specific targets.
Here,
the best defense is to very suspicious emails,
e.g., by phoning people to confirm their identity.
It's particularly important to do this when you are sent
things like password reset instructions.
Many IT departments send out messages that are indistinguishable from spearphishing attacks,
which just trains people to be victims.

\rulemajor{Rule 7: Erase before discarding.}

Moving files into the trash and then emptying it does not actually erase the data:
it just tells the computer that the space is available for re-use.
(This is why reporters and private investigators regularly go dumpster diving.)
The best way to address this problem is to encrypt the drive,
which is a quick setup option for all major operating systems these days.

Even with that,
you should act as if any device you throw away is going to fall into unfriendly hands.
Use a secure deletion tool like BleachBit on Linux or Windows
before selling, recycling, or discarding your hardware,
but keep in mind that this doesn't affect backups or files stored online
on sites like Dropbox.
And keep in mind that it is practically impossible to truly delete data from social networking sites:
in most cases,
their ``delete'' usually means ``don't show any more''
rather than ``erase all past record of''.

\rulemajor{Rule 8: Check your devices and accounts periodically.}

Google, Facebook, and others make money selling targeted advertising to you and data about you.
They have been forced to give users a semblance of control over personal data,
but they frequently change their terms of service in opaque ways.
Seemingly-innocuous information can give attackers valuable clues:
restaurant ``likes'' reveal where you were at specific times,
while funny stories about childhood birthday parties reveal likely answers to security questions.
Take advantage of what little protection they have been forced to allow you;
again,
if you get into the habit of doing this every time you do some other regular task,
you're more likely to do both.

Unfortunately,
even if you do this,
information may leak through others.
For example,
attackers can friend your friends in an attempt to get information about you,
such as the name of your first school.
And as bad as major social media sites are,
cell phone applications are often worse
(not even counting the ones that turn out to be government-sponsored spyware \cite{Schn2019}).
If a game wants access to your camera and address book,
you should probably find a different game to play.

Since social media is a fact of life for most of us,
you should check your settings periodically,
just as you would take your car in for an oil change.
(In fact,
the author does these things at the same time in order to remember both.)
Turn off everything you can
and then use a tracking blocker such as Ghostery to reduce information leakage.

Many experts recommend using separate devices or accounts for work and personal life,
but this is increasingly unrealistic.
Everyone checks their personal email from their work device eventually,
and everyone uses their personal phone for 2FA.
However,
you should consider getting a second phone for international travel:
the legalities around who can take your devices
and/or force you to unlock them are complicated and frequently overstepped,
so you should assume that anything on or connected to the devices you are traveling with
will be compromised.

\begin{quote}
  \emph{Never} plug a random USB drive into your device:
  it's like letting a complete stranger into your home unsupervised.
\end{quote}

\rulemajor{Rule 9: Fight back.}

Casual attackers may eventually get bored and move on,
but like all bullies,
they will also often revisit previous victims,
and even if they don't,
they are likely to pick new ones.
If you have been attacked:

\begin{enumerate}

\item
  \textbf{Find support.}
  Being targeted is frightening and wearying,
  particularly if you belong to one of the many groups that are targeted in real life
  as well as online.
  Let family, friends, and colleagues know what is happening so that they can support you.
  They may also be able to offer advice if they have been in similar situations.
  
\item
  \textbf{Document everything.}
  Save emails
  and take screenshots of sites like Facebook and Twitter
  (in case attackers delete or alter material).

\item
  \textbf{Do not engage directly.}
  Casual attackers are often seeking attention,
  so a direct response often encourages further attacks
  (and can draw attention from like-minded attackers).
  
\item
  \textbf{Report the attack.}
  Social media sites have done everything they can
  to avoid being accountable for facilitating online attacks,
  but companies and universities will usually take what steps they can
  once they know there is a problem.
  In the author's experience,
  they are most likely to do this if they believe that
  you might speak publicly about what has happened
  and thereby damage their reputation,
  so never agree to a non-disclosure agreement that would prevent you from doing so.
  
\end{enumerate}

\rulemajor{Rule 10: It's not all about you.}

Our final rule brings us full circle to the first one.
We don't just wear masks to prevent ourselves from becoming infected:
we also wear them so that we will not infect others.
Similarly,
if you do not take precautions with online security then you are putting others at risk.
Simple steps like putting passwords on PDFs that contain sensitive information
can go a long way to deter attackers,
in the same way that a sturdy-looking bike lock encourages would-be thieves to go after some other bike.
And if you \emph{are} compromised,
let those affected know as soon as you can.

But the only long-term way to improve everyone's online safety
is to pressure politicians to strengthen liability legislation
so that companies, universities, and other institutions have real incentives to take meaningful action.
Cars and drugs are as safe as they are because their manufacturers are liable for negligence and harm.
The sooner software companies and social media sites are liable as well,
the safer all of us will be.

\begin{quote}
  The author is grateful to Claire Bowen (Urban Institute),
  Leigh Honeywell (Tall Poppy),
  Mike Hoye (Mozilla),
  Scott Jackson (RStudio),
  Rick Johnson (RStudio),
  Cheng Soon Ong (CSIRO),
  and Sydney Young (EFF)
  for their feedback on this article.
\end{quote}

\nocite{*}
\bibliography{10-safety}

\section*{Appendix: VPNs and Tor}
  
A Virtual Private Network (VPN) connects your device to a server,
then has the server make connections to other machines on your behalf.
All messages between your device and the server are encrypted,
and the server can be managed by professional IT staff
in a jurisdiction with tight privacy laws
to increase your safety.
A specialized web browser called Tor routes messages randomly through a network of servers,
making traffic much harder to track.
Both of these reduce risk,
but neither eliminates it
if your device has been compromised,
if the VPN is compromised (or subpoenaed),
or if you log in to accounts over Tor
(thereby revealing your identity to those sites).

\end{document}
